% Section fragment; input from the main document.
\section{问题分析}
\label{sec:frontmatter-analysis}

\subsection{问题一的分析}

首先将一天划分为144个10分钟时段，把负载与光伏功率换算为时段电量。随后，以购电量、充放电量及储电量为变量，建立全天购电费用最小的线性规划，将供需平衡、储能递推、容量与功率限制及日循环条件纳入约束。求解后汇总指定时段的购电与充放电方案，核验能量平衡和设备边界，并与无储能方案比较购电费用。

\subsection{问题二的分析}

问题二按情景构造、计划决策与实际曲线回测三步展开。先依据月度差异、周周期及近期供需特征，筛选同季节历史完整日曲线；通过核加权赋予情景概率，以能量分数确定带宽。再建立两阶段随机线性规划：第一阶段确定各情景共用的日前购电量，第二阶段安排充放电、紧急购电与剩余电量弃置，最小化期望日费用。最后固定购电计划，在实测曲线上求解补救方案，将日末储电量传至次日，逐日计算费用；另对赋权方案和跨日终端价值作对照检验。

\subsection{问题三的分析}

首先利用历史数据预测负载，将整点光伏预报经保形插值与区间积分转为10分钟电量，再叠加同季节历史日的联合残差构造情景。在0:00求解条件两阶段随机规划，得到原计划；6:00、12:00和18:00到来时，冻结已执行时段，依据新预报和当前储电量重建情景，综合计划、调整及紧急购电费用与跨日终端价值，优化剩余计划，形成四阶段滚动策略。执行层采用情景模型预测控制（MPC），每10分钟重算并只执行当前动作。最后比较不同预报更新组合的实际结算费用，判断各时刻修正计划的必要性。

\subsection{问题四的分析}

先利用已知历史价格，按日内、周内和季节三个尺度构建电价预测，再抽取同一历史日的负载、光伏与电价残差，生成联合情景。随后将情景电价纳入购电目标与储能控制，分别求解日内不修改计划、按预报滚动修正计划两种方案，完成问题二、三的波动电价重算。最后按实际电价结算，比较两种策略的费用，并设置提前已知全天价格的对照，评价价格预测与信息更新的经济价值。

四个问题层层递进：问题一在确定条件下建立购电与储能的联合调度基线；问题二引入负载与光伏的不确定性，建立两阶段随机规划；问题三在此基础上加入日内预报更新，形成四阶段滚动与情景 MPC；问题四再把电价由固定扩展为实时波动。四个问题共享同一套能量平衡与储能递推约束，仅改变不确定性来源与决策时点，从而保证结果横向可比。
